\documentclass[11pt,a4paper]{article}

\usepackage[utf8]{inputenc}
\usepackage[T1]{fontenc}
\usepackage[ngerman]{babel}

\usepackage{amsmath,amssymb,amsthm,mathtools}
\usepackage[a4paper,margin=2.5cm]{geometry}
\usepackage{lmodern}
\usepackage{microtype}
\usepackage{hyperref}
\usepackage{enumitem}

\title{Eine ultrakurze lokale Strukturtheorie des EABC-Gitters modulo \(12\)}
\author{Kollaborative Ausarbeitung}
\date{\today}

\newtheorem{satz}{Satz}
\newtheorem{lemma}[satz]{Lemma}
\theoremstyle{definition}
\newtheorem{definition}[satz]{Definition}
\theoremstyle{remark}
\newtheorem{bemerkung}[satz]{Bemerkung}

\newcommand{\N}{\mathbb{N}}
\newcommand{\ind}{\mathbf{1}}

\begin{document}
	\maketitle
	
	\begin{abstract}
		Wir formulieren eine ultrakurze lokale Strukturtheorie des modulo-\(12\)-Gitters
		\[
		Q_n=(12n-11,\,12n-7,\,12n-5,\,12n-1).
		\]
		Dabei werden Primzahlen \(>3\), Primzahlzwillinge, Primzahldrillinge,
		blockübergreifende Primzahlvierlinge und dichte Primzahlfünflinge in einem einheitlichen
		Schema beschrieben.
	\end{abstract}
	
	\begin{definition}[EABC-Gitter]
		Für \(n\in\N\) setzen wir
		\[
		E_n:=12n-11,\qquad
		A_n:=12n-7,\qquad
		B_n:=12n-5,\qquad
		C_n:=12n-1,
		\]
		und nennen
		\[
		Q_n:=(E_n,A_n,B_n,C_n)
		\]
		den \(n\)-ten EABC-Block. Ferner definieren wir die beiden Zwillingskanäle
		\[
		T_n^{\mathrm{int}}:=(A_n,B_n)=(12n-7,\,12n-5),
		\]
		\[
		T_n^{\mathrm{ext}}:=(C_n,E_{n+1})=(12n-1,\,12n+1),
		\]
		sowie die Zwillingssignatur
		\[
		\tau(Q_n):=
		\bigl(
		\ind_{T_n^{\mathrm{int}}\text{ ist Primzahlzwilling}},
		\ind_{T_n^{\mathrm{ext}}\text{ ist Primzahlzwilling}}
		\bigr)\in\{0,1\}^2.
		\]
	\end{definition}
	
	\begin{lemma}[Modulo-\(12\)-Zerlegung und Zwillingskanäle]
		\leavevmode
		\begin{enumerate}[label=\textup{(\roman*)}]
			\item Jede Primzahl \(p>3\) liegt eindeutig in genau einer der vier Familien
			\[
			E,\ A,\ B,\ C,
			\]
			das heißt in genau einer der vier Restklassen
			\[
			1,5,7,11 \pmod{12}.
			\]
			\item Jeder Primzahlzwilling \((p,p+2)\) mit \(p>3\) liegt eindeutig in genau einem der beiden Muster
			\[
			(12n-7,12n-5)
			\qquad\text{oder}\qquad
			(12n-1,12n+1).
			\]
		\end{enumerate}
	\end{lemma}
	
	\begin{proof}
		Zu \textup{(i)}: Für \(p>3\) ist \(p\) weder durch \(2\) noch durch \(3\) teilbar. Daher gilt
		\[
		p\equiv 1,5,7,\text{ oder }11 \pmod{12}.
		\]
		Diese vier Restklassen werden durch \(E_n,A_n,B_n,C_n\) disjunkt und vollständig parametrisiert.
		
		Zu \textup{(ii)}: Für einen Primzahlzwilling \((p,p+2)\) mit \(p>3\) sind modulo \(12\) nur die Muster
		\[
		(5,7)\quad\text{oder}\quad(11,1)
		\]
		möglich. Diese entsprechen genau den Paaren
		\[
		(A_n,B_n)\quad\text{bzw.}\quad(C_n,E_{n+1}).
		\]
	\end{proof}
	
	\begin{lemma}[Drillinge, Vierlinge und Fünflinge]
		\leavevmode
		\begin{enumerate}[label=\textup{(\roman*)}]
			\item Jeder dichte Primzahldrilling enthält die Primzahl \(3\); daher sind die einzigen dichten
			Primzahldrillinge
			\[
			(3,5,7)
			\qquad\text{und}\qquad
			(3,7,11).
			\]
			\item Es gilt
			\[
			\tau(Q_n)=(1,1)
			\]
			genau dann, wenn
			\[
			(12n-7,\,12n-5,\,12n-1,\,12n+1)
			\]
			ein blockübergreifender Primzahlvierling ist.
			\item Jeder dichte Primzahlfünfling mit \(p>3\) besitzt genau eine der beiden Formen
			\[
			(A_n,B_n,C_n,E_{n+1},A_{n+1})
			\]
			oder
			\[
			(C_n,E_{n+1},A_{n+1},B_{n+1},C_{n+1}),
			\]
			und ist damit eine orientierte Erweiterung einer Doppelstelle.
		\end{enumerate}
	\end{lemma}
	
	\begin{proof}
		Zu \textup{(i)}: In einem dichten Drillingsmuster
		\[
		(p,p+2,p+6)\quad\text{oder}\quad(p,p+4,p+6)
		\]
		ist stets eine Zahl durch \(3\) teilbar. Ist sie prim, so muss sie gleich \(3\) sein. Das liefert
		genau die beiden Fälle \((3,5,7)\) und \((3,7,11)\).
		
		Zu \textup{(ii)}: Die Bedingung \(\tau(Q_n)=(1,1)\) bedeutet, dass sowohl
		\[
		(12n-7,12n-5)
		\]
		als auch
		\[
		(12n-1,12n+1)
		\]
		Primzahlzwillinge sind. Das ist äquivalent zur Primheit aller vier Zahlen
		\[
		12n-7,\ 12n-5,\ 12n-1,\ 12n+1,
		\]
		also zu einem Vierling der Form \(p,p+2,p+6,p+8\).
		
		Zu \textup{(iii)}: Für einen dichten Fünfling
		\[
		(p,p+2,p+6,p+8,p+12)
		\]
		sind als Startreste modulo \(12\) nur \(5\) oder \(11\) möglich. Daraus ergeben sich unmittelbar die
		beiden Formen
		\[
		(A_n,B_n,C_n,E_{n+1},A_{n+1})
		\quad\text{und}\quad
		(C_n,E_{n+1},A_{n+1},B_{n+1},C_{n+1}).
		\]
		In beiden Fällen enthalten die ersten vier Einträge eine Doppelstelle.
	\end{proof}
	
	\begin{satz}[Lokale Vollständigkeit des EABC-Gitters]
		Das EABC-Gitter
		\[
		Q_n=(12n-11,12n-7,12n-5,12n-1)
		\]
		liefert eine vollständige lokale Beschreibung der kleinen Primzahlkonstellationen:
		
		\begin{enumerate}[label=\textup{(\roman*)}]
			\item Primzahlen \(>3\) zerfallen eindeutig in die vier Familien \(E,A,B,C\).
			\item Primzahlzwillinge \(>3\) laufen genau über die beiden Kanäle
			\[
			(A_n,B_n)\quad\text{und}\quad(C_n,E_{n+1}).
			\]
			\item Dichte Primzahldrillinge sind genau die beiden singulären \(3\)-verankerten Fälle
			\[
			(3,5,7),\qquad (3,7,11).
			\]
			\item Doppelstellen \(\tau(Q_n)=(1,1)\) sind genau blockübergreifende Primzahlvierlinge
			\[
			(12n-7,12n-5,12n-1,12n+1).
			\]
			\item Dichte Primzahlfünflinge sind genau die orientierten Erweiterungen solcher Doppelstellen,
			also genau die beiden Muster
			\[
			(A_n,B_n,C_n,E_{n+1},A_{n+1})
			\quad\text{und}\quad
			(C_n,E_{n+1},A_{n+1},B_{n+1},C_{n+1}).
			\]
		\end{enumerate}
		
		Damit ist die gesamte lokale Zwillings-, Drillings-, Vierlings- und Fünflingsstruktur
		des Primzahlgitters modulo \(12\) vollständig beschrieben.
	\end{satz}
	
	\begin{proof}
		Die fünf Aussagen folgen unmittelbar aus den beiden Lemmata.
	\end{proof}
	
	\begin{bemerkung}
		In dieser Sicht ist der Vierling die gleichzeitige Aktivierung beider Zwillingskanäle,
		während der Fünfling die erste gerichtete Erweiterung einer solchen Doppelstelle darstellt.
	\end{bemerkung}
	
\end{document}